\section*{Impact Statement}
This work aims to make explanations of language-model internals more reliable, by showing that which cache channel appears to carry a value, and whether an edit at a context token reaches the answer at all, depend on the prompt under which they are measured. More reliable component-level explanations could help researchers audit and debug models, and could reduce the risk that a mechanism found under one evaluation format is trusted to explain behaviour under another. Our results also caution against over-reading interpretability findings: an attribution or an edit validated under one prompt format, or inferred from the model's own read, should not on its own certify behaviour under others; an edit that works in free form can fail silently once the candidates are listed. Methods for locating and editing what a model reads from its context could also be used to steer models in unwanted ways; our experiments use public models, templated stories, public reading-comprehension passages, edits built from public material, a released remap and a released state-token patch.
